\section{Related work and contribution boundary}\label{related-work-and-contribution-boundary}

POEM \cite[3.2]{bar2024poem} already exhibits scalar Gaussian entropy collapse; imbalance-driven deterioration and stabilization are established \cite{weerasinghe2026}. Sequential risk monitoring \cite{schirmer2025} and unlabeled action identifiability \cite{jhawar2026} have different observation timing and targets. Our full risk function is identical across environments, while score marginals differ locally and the target is an independent future path probability. Merely changing assumptions is not the contribution; the quantitative shared-risk prediction result is.

The random-seed amplification mechanism is classical unstable-exit theory \cite{bakhtin2008}. Experiment comparison and bounded-loss transfer \cite{nussbaum2004}, influence-aligned least-favorable directions \cite{vanderlaan2016}, and posterior medians under absolute loss are classical tools. The new model-specific objects are the legal common-risk family, the uniform real-SGD exit interface and the resulting information/precision characterization. The focused literature table records the comparisons and the limits of its source audit; it is not an exhaustive priority claim.

\section{Limitations and discussion}\label{limitations-and-discussion}

The initial results assume a known baseline; the finite-baseline extension estimates its root and scales within a fixed regular density neighborhood, with L/n tending to a positive limit or infinity. Neither setting covers arbitrary target distributions. The envelope is uniformly bounded, perturbations are of order \(n^{-1/2}\), the seed window is fixed, and the least-favorable subfamily must fit in that envelope. Step size and risk tolerance stay fixed and positive. Momentum, schedules, temporal dependence, batch reuse, online trajectory information and feature learning require separate analysis.

The minimax transfer is asymptotic, with no numerically usable bound on experiment deficiency or the exit remainder. The Gaussian intervals certify a limiting decision value; the confirmation intervals certify finite-grid comparisons in the stated simulation experiment. Neither certifies continuous-window finite-n minimax optimality. The optimized-rule confirmation does not test nuisance directions; the separate finite-direction check tests the raw moment interface, and a known risk function is an analytical control, not an unlabeled risk-estimation method.

The original Gaussian geometry and exit-time results remain in the supporting Gaussian-model chapter and the Gaussian-proof chapter. Earlier frozen-feature studies and failed stopping tests remain in the supporting experimental chapter with their individual limitations. They motivate calibration sensitivity but do not enlarge the scope to general neural-network adaptation. The present work supplies a prospective probability-prediction theory and a model-specific validation; it does not supply a general unlabeled stopping algorithm.

The visual confirmation conditions on a known official-test empirical baseline and independent draws with replacement from fixed image pools. Its environment-specific weighted label risk differs across tilts, so it is a mechanism and prediction validation rather than a test of the common-risk construction. Automobile/truck does not have a positive simultaneous improvement interval over clipping; no pair was excluded or given extra confirmation budget. The natural-corruption primary comparison remains unresolved; finite-baseline guarantees remain experimentally untested.

The scalar update adapts one threshold on a frozen projection with actual entropy gradients. It retains batch randomness, local instability and risk exit, but is not an established reduction of BN-affine Tent. Exit side describes threshold bias rather than inevitable class collapse. The full-network check in Section~\ref{full-network-bn-development-checks} supplies an empirical boundary: a highly reproducible mean direction coexists with poorly transferring residual axes and improving risk. It therefore supports neither a dominant one-dimensional stochastic mode nor the declared risk-event forecast. This applicability result distinguishes evidence for the scalar mechanism under controlled conditions from evidence needed for its transfer; it does not refute a theorem whose unstable-root assumptions have not been established in the tested network.
